\subsection{Empirical implications}\label{app:empirical}

This subsection states and proves the results that Section~\ref{sec:empirical} summarizes. The incumbent law \(F\)
is any continuous law, the class laws satisfy the class ranking \eqref{eq:classes}, and the prior \(\pi\in(0,1)\) is
arbitrary; as in Lemma~\ref{lem:prices}, \(\DT>0\) and \(\gamma=g_H-g_L>0\). Each result names the trading pattern
and the tie rules it uses. Every statement concerns one target in one public state, with its prior, slot and
incumbent law fixed, and every probability, mean and covariance is taken over the order opportunity of that state.
Write \(M=M(\pi)\) and \(m=m(\pi)\); \(\bar f(x)=\Pr(X=x)\) for the frequency of the order \(x\) per order
opportunity; \(v=\Pr(X\ne0)\) for the volume; \(w_F(\mu)=t_L-t_0+\DT\mu\) for the entry gain at the belief \(\mu\)
against \(F\) (Appendix~\ref{app:beliefs}); and \(\chi_F=\1\{B_F(M)\ge c_H\}\). The entry gain rises in \(\mu\) with
slope \(\DT\) and is positive at every interior belief. By Lemma~\ref{lem:prices}(d), entry after an order is \(\rho\)
when its posterior lies below \(\tau\) and one when it lies above; at the posterior \(\tau\) it is one when an
indifferent high-cost challenger prepares, and any value in \([\rho,1]\) when an indifferent high-cost challenger
may prepare with any probability.

The purchase impact is \(I_+=P(1)-P(0)\), the sale impact is \(I_-=P(0)-P(-1)\), and the price gap
\(P(1)-P(-1)=I_++I_-\) is the difference between the conditional execution prices of a unit purchase and a unit
sale. The model has no separate quote stage: each price is the competitive value conditional on the order. The
weighted impact difference is \(\bar f(1)I_+-\bar f(-1)I_-\).

Under full two-sided trade, \(\sigma_H=\delta_1\) and \(\sigma_L=\delta_{-1}\), Lemma~\ref{lem:beliefs}(d) gives
the posteriors \((m,\pi,M)\) after the orders \((-1,0,1)\), and
\begin{equation}\label{eq:a8-full}
\begin{gathered}
\bar f(1)=b+\pi\lambda,\qquad \bar f(0)=1-2b-\lambda,\qquad \bar f(-1)=b+(1-\pi)\lambda,\\
v=2b+\lambda,\qquad \E[X]=\lambda(2\pi-1).
\end{gathered}
\end{equation}
Subtracting the prior from each posterior gives
\begin{equation}\label{eq:a8-balance}
\bar f(1)\,(M-\pi)=\bar f(-1)\,(\pi-m)=\pi(1-\pi)\lambda .
\end{equation}
The informed share of nonzero orders is \(\vartheta=\lambda/v=u/(2+u)\) at every prior. The posterior gap is
\(M-m=\pi(1-\pi)\lambda v/[\bar f(1)\bar f(-1)]\le\vartheta\), with equality exactly at the prior \(1/2\), where \(M+m=1\)
and \(\vartheta=2M-1=M-m\). When \(\pi<\tau\), entry after a sale and after no order is \(\rho\) under both challenger
rules.

\begin{proof}[Proof of the full-trade identities]
Under full two-sided trade, class \(H\) places the orders \((-1,0,1)\) with probabilities \((b,\,1-2b-\lambda,\,b+\lambda)\)
and class \(L\) with probabilities \((b+\lambda,\,1-2b-\lambda,\,b)\). Averaging with the weights \(\pi\) and \(1-\pi\)
gives \eqref{eq:a8-full}, and Lemma~\ref{lem:beliefs}(d) gives the posteriors. Bayes' rule gives
\(\bar f(1)\,M=\pi(b+\lambda)\) and \(\bar f(-1)\,m=\pi b\). Subtracting \(\pi\bar f(1)\) from the first and
\(\pi\bar f(-1)\) from the second gives
\(\bar f(1)(M-\pi)=\pi[b+\lambda-\bar f(1)]=\pi(1-\pi)\lambda\) and
\(\bar f(-1)(\pi-m)=\pi[\bar f(-1)-b]=\pi(1-\pi)\lambda\), which is \eqref{eq:a8-balance}. Adding \(M-\pi\) and
\(\pi-m\), with \(\bar f(1)+\bar f(-1)=v\), gives the posterior gap. With \(u=\lambda/b\), \(\lambda/v=\lambda/(2b+\lambda)=u/(2+u)\).
The posterior gap is \(\vartheta\) times \(\pi(1-\pi)v^2/[\bar f(1)\bar f(-1)]\), and
\[
\bar f(1)\,\bar f(-1)-\pi(1-\pi)v^2=b^2+b\lambda+\pi(1-\pi)\lambda^2-\pi(1-\pi)(2b+\lambda)^2=b(b+\lambda)(1-2\pi)^2\ge0,
\]
with equality exactly at the prior \(1/2\), because \(b>0\). At the prior \(1/2\), \(M=(b+\lambda)/v\) and \(m=b/v\), so
\(M+m=1\) and \(2M-1=\lambda/v\). The entry gain is \(w_F(\mu)=w_L+\DT\mu\ge\DT\mu\), since \(w_L\ge0\) (Appendix~\ref{app:beliefs}),
and is therefore positive at every \(\mu>0\). When \(\pi<\tau\), the posteriors after a sale and after no order are
\(m<\pi<\tau\), and Lemma~\ref{lem:prices}(d) gives entry \(\rho\) after both, with no posterior at \(\tau\). These steps use
parts (a), (b) and (d) of Lemma~\ref{lem:beliefs}, which hold for every \(\beta<1/2\).
\end{proof}

\begin{proposition}[The price gap and the weighted test]\label{prop:spread}
\leavevmode
\begin{enumerate}[label=(\roman*),leftmargin=2em]
\item \emph{Every trading pattern.} In every equilibrium,
\[
\bar f(1)I_+-\bar f(-1)I_-=\E[P(X)]-P(0)=e(0)\,\DT\,[\pi-\mu(0)]+\sum_x\bar f(x)\,[e(x)-e(0)]\,w_F(\mu(x)).
\]
\item \emph{The weighted test.} Under full two-sided trade,
\[
\bar f(1)I_+-\bar f(-1)I_-=\bar f(1)\,[e(1)-e(0)]\,w_F(M)-\bar f(-1)\,[e(0)-e(-1)]\,w_F(m).
\]
If also \(\pi<\tau\), then
\[
I_+-\frac{\bar f(-1)}{\bar f(1)}\,I_-=[e(1)-\rho]\,w_F(M),
\]
which is positive exactly when a purchase recruits, that is, exactly when \(e(1)>\rho\), under both challenger rules.
\item \emph{Signs.} Let full two-sided trade hold, and let an indifferent high-cost challenger prepare. Entry after
the orders \((-1,0,1)\) and the weighted impact difference are:
\begin{center}
\begin{tabular}{@{}lll@{}}
\toprule
Preparation threshold & Entry after \((-1,0,1)\) & \(\bar f(1)I_+-\bar f(-1)I_-\)\\
\midrule
\(\tau\le m\) & \((1,1,1)\) & \(0\)\\
\(m<\tau\le\pi\) & \((\rho,1,1)\) & \(-(1-\rho)\,\bar f(-1)\,w_F(m)<0\)\\
\(\pi<\tau\le M\) & \((\rho,\rho,1)\) & \((1-\rho)\,\bar f(1)\,w_F(M)>0\)\\
\(M<\tau\) & \((\rho,\rho,\rho)\) & \(0\)\\
\bottomrule
\end{tabular}
\end{center}
When an indifferent high-cost challenger may prepare with any probability, the table holds for
\(\tau\notin\{m,\pi,M\}\). At \(\tau=M\) the weighted impact difference is \([e(1)-\rho]\,\bar f(1)\,w_F(M)\ge0\),
positive exactly when \(e(1)>\rho\); at \(\tau=m\) it is \(-[1-e(-1)]\,\bar f(-1)\,w_F(m)\le0\), negative exactly when
\(e(-1)<1\). At \(\tau=\pi\), full two-sided trade is an equilibrium with one entry \(e(0)\in[\rho,1]\) after no
order exactly when it is an equilibrium with every such entry, and as \(e(0)\) runs over \([\rho,1]\) the weighted
impact difference \(\bar f(1)\,[1-e(0)]\,w_F(M)-\bar f(-1)\,[e(0)-\rho]\,w_F(m)\) takes every value from
\(-(1-\rho)\,\bar f(-1)\,w_F(m)\) to \((1-\rho)\,\bar f(1)\,w_F(M)\).
\item \emph{The price gap.} Under full two-sided trade with \(\pi<\tau\),
\[
P(1)-P(-1)=\rho\DT\,(M-m)+[e(1)-\rho]\,w_F(M),\qquad \rho\DT\,(M-m)=\frac{v}{\bar f(1)}\,I_-,
\]
a class component and an entry component. If the trading margin \(\rho\kappa(\pi)\DT-k\) is positive, with
\(\kappa(\pi)=\min\{m,1-M\}\), every equilibrium has full two-sided trade, under both investor rules, whether or not
the challenger sees the price. If also \(\pi<\tau\) and the challenger does not see the price, its belief stays at
\(\pi\), entry is \(\rho\) after every order, the price gap is the class component \(\rho\DT\,(M-m)\), and the
weighted impact difference is zero.
\item \emph{The unweighted comparison.} Under full two-sided trade with \(\pi<\tau\),
\[
I_+-I_-=\rho\DT\,\frac{\pi(1-\pi)\lambda^2(1-2\pi)}{\bar f(1)\,\bar f(-1)}+[e(1)-\rho]\,w_F(M).
\]
If \(\pi\ge1/2\), \(\rho\le1/2\) and an indifferent high-cost challenger prepares, then under full two-sided trade
\(I_+>I_-\) exactly when \(\pi<\tau\le M\), that is, exactly when a purchase raises entry above its level after no
order.
\item \emph{Strength.} Let \(F_r\) be a path along which \(\DT\) rises strictly, hold the prior and the slot fixed,
and write \(w_r=w_{F_r}\) and \(\chi_r=\chi_{F_r}\).
\begin{enumerate}[label=(\alph*),leftmargin=2em]
\item For every belief \(\mu\),
\[
w_r(\mu)=pF_r(p)+\int_p^\infty\bigl[1-(1-\mu)G_L(y)-\mu G_H(y)\bigr]\bigl[1-F_r(y)\bigr]\,dy:
\]
the value of rescuing the sale when the incumbent misses the reserve, plus the value of competition above the
reserve. From \(r\) to \(r'>r\), the first term falls by \(p\,[F_r(p)-F_{r'}(p)]\ge0\), and the second rises by
\(\int_p^\infty[1-(1-\mu)G_L(y)-\mu G_H(y)]\,[F_r(y)-F_{r'}(y)]\,dy\ge\mu\,[\DT(r')-\DT(r)]\). The entry gain
\(w_r(\mu)\) rises from \(r\) to \(r'\) exactly when the rise of the second term exceeds the fall of the first.
\item If the probability \(F_r(p)\) that the incumbent misses the reserve is the same at every strength of the path,
\(w_r(\mu)\) rises strictly in \(r\) for every \(\mu\in(0,1]\). If \(F_r\) is uniform on \([0,r]\) with \(r>p\), then
\[
\frac{d}{dr}\,w_r(\mu)=\frac{\E_\mu\bigl[\min\{\theta,r\}^2\bigr]-3p^2}{2r^2},
\]
where \(\E_\mu\) takes \(\theta\) from \((1-\mu)G_L+\mu G_H\). So \(w_r(\mu)\) rises strictly in \(r\) for every
\(\mu\) when \(p<\min\{\underline\theta,\underline r\}/\sqrt3\), where \(\underline\theta\) is the lowest challenger
value and \(\underline r\) the lowest strength on the path.
\item Let an indifferent high-cost challenger prepare, and let full two-sided trade hold at every strength of an
interval on which \(\pi<\tau(r)\) and \(w_r(M)\) rises strictly. Where \(\chi_r=1\), both components of the price gap,
both impacts, the price gap and the weighted impact difference rise strictly with \(r\). Where \(\chi_r=0\), the
entry component and the weighted impact difference are zero, and the class component, both impacts and the price
gap rise strictly with \(r\).
\item The indicator \(\chi_r\) is nonincreasing in \(r\). Let also \(F_r(y)\) be continuous in \(r\) at every \(y\);
then the strengths with \(\chi_r=1\) form a closed subset of the path. Let \(r_C\) be a strength with
\(\chi_{r_C}=1\) and \(\chi_r=0\) for every \(r>r_C\), and let full two-sided trade and \(\pi<\tau(r)\) hold on a
neighborhood of \(r_C\). When an indifferent
high-cost challenger prepares, \(P(1)\), \(I_+\) and the price gap are continuous from the left at \(r_C\), and their
right limits lie below their values at \(r_C\) by \((1-\rho)\,w_{r_C}(M)>0\); the weighted impact difference falls
by \((1-\rho)\,\bar f(1)\,w_{r_C}(M)\) in the same way; \(P(-1)\), \(P(0)\) and \(I_-\) are continuous at \(r_C\).
When an indifferent investor may mix over its best orders and an indifferent high-cost challenger may prepare with
any probability, entry after a purchase at \(r_C\) can take
any value in \([\rho,1]\), and the price gap at \(r_C\) any value between its two one-sided limits.
\end{enumerate}
\end{enumerate}
\end{proposition}

\begin{proof}[Proof of Proposition~\ref{prop:spread}]
By Lemma~\ref{lem:prices}(b), \(P(x)=t_0+e(x)\,w_F(\mu(x))\) after every order, in every equilibrium and under both
challenger rules.
\begin{enumerate}[label=(\roman*),leftmargin=2em]
\item Since \(P(1)-P(0)=I_+\) and \(P(-1)-P(0)=-I_-\),
\(\E[P(X)]-P(0)=\sum_x\bar f(x)\,[P(x)-P(0)]=\bar f(1)I_+-\bar f(-1)I_-\). Every order has positive probability by
Lemma~\ref{lem:beliefs}(a), so Bayes' rule gives \(\sum_x\bar f(x)\,\mu(x)=\pi\), and since \(w_F\) is affine,
\(\sum_x\bar f(x)\,w_F(\mu(x))=w_F(\pi)\). Hence
\[
\begin{aligned}
\E[P(X)]-P(0)&=\sum_x\bar f(x)\,e(x)\,w_F(\mu(x))-e(0)\,w_F(\mu(0))\\
&=\sum_x\bar f(x)\,[e(x)-e(0)]\,w_F(\mu(x))+e(0)\,[w_F(\pi)-w_F(\mu(0))],
\end{aligned}
\]
and \(w_F(\pi)-w_F(\mu(0))=\DT\,[\pi-\mu(0)]\).
\item Under full two-sided trade \(\mu(0)=\pi\), so the first term of (i) is zero, and the posteriors after the orders
\(1\) and \(-1\) are \(M\) and \(m\). The sum in (i) then has the two terms of the display. If \(\pi<\tau\), entry after
a sale and after no order is \(\rho\) by the standing paragraph, so the display equals \(\bar f(1)\,[e(1)-\rho]\,w_F(M)\);
dividing by \(\bar f(1)>0\) gives the second display. Since \(w_F(M)>0\) and \(e(1)\ge\rho\), it is positive exactly when
\(e(1)>\rho\). No step uses the challenger's rule at a posterior equal to \(\tau\).
\item Under full two-sided trade \(m<\pi<M\). When an indifferent high-cost challenger prepares,
Lemma~\ref{lem:prices}(d) gives \(e(x)=1\) when \(\mu(x)\ge\tau\) and \(e(x)=\rho\) otherwise, which is the entry column.
The display of (ii) gives the last column, and \(w_F(m),w_F(M)>0\) give the signs. When an indifferent high-cost
challenger may prepare with any probability, Lemma~\ref{lem:prices}(d) differs only after an order whose posterior
equals \(\tau\), so the table holds for \(\tau\notin\{m,\pi,M\}\). At \(\tau=M\), \(e(-1)=e(0)=\rho\) and
\(e(1)\in[\rho,1]\); at \(\tau=m\), \(e(0)=e(1)=1\) and \(e(-1)\in[\rho,1]\); the display of (ii) gives the two values.
At \(\tau=\pi\), \(e(-1)=\rho\), \(e(1)=1\), and \(e(0)\in[\rho,1]\) is the challenger's choice at the posterior \(\tau\).
The posteriors do not depend on entry, and by Lemma~\ref{lem:trading}(d) entry after the zero order enters no payoff
of the investor, so the investor's conditions are the same for every \(e(0)\). Entry is nondecreasing in the
posterior for every \(e(0)\in[\rho,1]\), so the paragraph on constructed tuples in Appendix~\ref{app:beliefs} supplies
the prices, the preparation rule and the price beliefs, and (E3) holds at the posterior \(\tau\). Full two-sided trade
is therefore an equilibrium with one such \(e(0)\) exactly when it is an equilibrium with every such \(e(0)\). By
(ii), the weighted impact difference is \(\bar f(1)\,[1-e(0)]\,w_F(M)-\bar f(-1)\,[e(0)-\rho]\,w_F(m)\), which is affine
and strictly decreasing in \(e(0)\), with the value \((1-\rho)\,\bar f(1)\,w_F(M)\) at \(e(0)=\rho\) and
\(-(1-\rho)\,\bar f(-1)\,w_F(m)\) at \(e(0)=1\).
\item With \(\pi<\tau\), \(e(-1)=e(0)=\rho\), so
\[
P(1)-P(-1)=e(1)\,w_F(M)-\rho\,w_F(m)=\rho\DT\,(M-m)+[e(1)-\rho]\,w_F(M).
\]
Also \(I_-=\rho\,[w_F(\pi)-w_F(m)]=\rho\DT\,(\pi-m)\). By \eqref{eq:a8-balance},
\(\pi-m=\pi(1-\pi)\lambda/\bar f(-1)\), and by the standing paragraph
\(M-m=\pi(1-\pi)\lambda v/[\bar f(1)\bar f(-1)]\), so \(M-m=(\pi-m)\,v/\bar f(1)\), which gives the second identity.

In the economy where the challenger sees the price, Lemma~\ref{lem:trading}(c) bounds \(K_H\) and \(K_L\) below by
\(\rho\kappa(\pi)\DT\) in every candidate. Where the challenger does not see the price, its belief is \(\pi\) after
every order, so it prepares with one probability \(e\in[\rho,1]\) after every order. Preparation is then independent
of the class and of the order, the competitive price is \(P(x)=t_0+e\,w_F(\mu(x))\) by the argument of
Lemma~\ref{lem:prices}(b), and the payoffs of Lemma~\ref{lem:trading} hold with \(e(x)=e\). Its bounds in part (c)
use only \(e\ge\rho\) and the posterior bounds of Lemma~\ref{lem:beliefs}(b), so they hold as well. In both economies,
\(\rho\kappa(\pi)\DT>k\) gives \(K_H>k\) and \(K_L>k\) in every candidate, and Lemma~\ref{lem:trading}(a) and (b) make
class \(H\) buy and class \(L\) sell with probability one, with no indifference, under both investor rules.

Let also \(\pi<\tau\), with a hidden price. Gross profit at the belief \(\pi\) is \(B_F(\pi)=c_H+\gamma(\pi-\tau)<c_H\),
so the high-cost challenger abstains, and entry is \(\rho\) after every order. The orders are full, the posteriors are
\((m,\pi,M)\), and \(P(x)=t_0+\rho\,w_F(\mu(x))\). The price gap is \(\rho\,[w_F(M)-w_F(m)]=\rho\DT\,(M-m)\). The
computation of (i) and (ii) uses only this price formula, Bayes' rule and full two-sided trade, so with \(e\equiv\rho\)
it gives a zero weighted impact difference.
\item With \(\pi<\tau\), \(I_+=e(1)\,w_F(M)-\rho\,w_F(\pi)=\rho\DT\,(M-\pi)+[e(1)-\rho]\,w_F(M)\) and
\(I_-=\rho\DT\,(\pi-m)\). By \eqref{eq:a8-balance},
\[
(M-\pi)-(\pi-m)=\pi(1-\pi)\lambda\,\frac{\bar f(-1)-\bar f(1)}{\bar f(1)\,\bar f(-1)},
\]
and \(\bar f(-1)-\bar f(1)=\lambda(1-2\pi)\) by \eqref{eq:a8-full}. This gives the formula.

Let \(\pi\ge1/2\), \(\rho\le1/2\), and let an indifferent high-cost challenger prepare. Then \(\bar f(1)\ge\bar f(-1)\),
and \eqref{eq:a8-balance} gives \(M-\pi\le\pi-m\). Take the rows of the table in (iii).
\begin{itemize}[leftmargin=2em]
\item \(\pi<\tau\le M\), entry \((\rho,\rho,1)\). The formula with \(e(1)=1\) and \(w_F(M)=w_L+\DT M\) gives
\(I_+-I_-=(1-\rho)\,w_L+\DT\,(M-2\rho\pi+\rho m)\). Since \(\rho\le1/2\), \(M-2\rho\pi\ge M-\pi>0\), and \(w_L\ge0\), so
\(I_+>I_-\).
\item \(M<\tau\), entry \((\rho,\rho,\rho)\). The formula gives \(I_+-I_-=\rho\DT\,[(M-\pi)-(\pi-m)]\le0\).
\item \(\tau\le m\), entry \((1,1,1)\). Then \(I_+=\DT\,(M-\pi)\) and \(I_-=\DT\,(\pi-m)\), so \(I_+\le I_-\).
\item \(m<\tau\le\pi\), entry \((\rho,1,1)\). Then \(I_+=\DT\,(M-\pi)\) and
\(I_-=w_F(\pi)-\rho\,w_F(m)=\rho\DT\,(\pi-m)+(1-\rho)\,w_F(\pi)\). Hence
\(I_+-I_-\le(1-\rho)\,\DT\,(\pi-m)-(1-\rho)\,w_F(\pi)<0\), because \(w_F(\pi)\ge\DT\pi>\DT\,(\pi-m)\).
\end{itemize}
So \(I_+>I_-\) exactly in the row \(\pi<\tau\le M\), the only row in which \(e(1)>e(0)\).
\item The prior and the slot are fixed, so \(M\), \(m\), \(\bar f(x)\) and \(v\) do not depend on \(r\).
\begin{enumerate}[label=(\alph*),leftmargin=2em]
\item Under the class ranking \(G_j(p^-)=0\), and Appendix~\ref{app:payoffs} gives
\(w_j=t_j-t_0=pF(p)+\int_p^\infty[1-G_j(y)][1-F(y)]\,dy\). Averaging with the weights \(1-\mu\) and \(\mu\) gives the
display. Along the path \(F_{r'}\le F_r\), which gives the change of the first term. The integrand of the change of the
second term carries \(1-(1-\mu)G_L-\mu G_H=1-G_L+\mu D\ge\mu D\), and Lemma~\ref{lem:twonumbers} gives
\(\int_p^\infty[F_r(y)-F_{r'}(y)]\,D(y)\,dy=\DT(r')-\DT(r)\). The last sentence compares the two changes.
\item If \(F_r(p)\) does not depend on \(r\), the first term does not change, and by (a) the second rises by at least
\(\mu\,[\DT(r')-\DT(r)]>0\) for \(\mu>0\). Let \(F_r\) be uniform on \([0,r]\) with \(r>p\), and write
\(\bar G_\mu(y)=1-(1-\mu)G_L(y)-\mu G_H(y)\), the probability that \(\theta>y\) under \((1-\mu)G_L+\mu G_H\). Then
\[
w_r(\mu)=\frac{p^2}{r}+\frac1r\int_p^r(r-y)\,\bar G_\mu(y)\,dy .
\]
The function \(r\mapsto\int_p^r(r-y)\bar G_\mu(y)\,dy=\int_p^r\int_p^z\bar G_\mu(y)\,dy\,dz\) has the continuous derivative
\(\int_p^r\bar G_\mu(y)\,dy\), so \(w_r(\mu)\) is differentiable in \(r\), with
\[
\frac{d}{dr}\,w_r(\mu)=-\frac{p^2}{r^2}+\frac{r\int_p^r\bar G_\mu(y)\,dy-\int_p^r(r-y)\bar G_\mu(y)\,dy}{r^2}
=\frac{\int_p^r y\bar G_\mu(y)\,dy-p^2}{r^2}.
\]
Every challenger value is at least \(p\), so \(\bar G_\mu=1\) below \(p\), and
\(\E_\mu[\min\{\theta,r\}^2]=\int_0^r2y\,\bar G_\mu(y)\,dy=p^2+2\int_p^ry\bar G_\mu(y)\,dy\). This gives the formula. On the path
\(r\ge\underline r\), and every value satisfies \(\theta\ge\underline\theta\), so
\(\min\{\theta,r\}\ge\min\{\underline\theta,\underline r\}>\sqrt3\,p\), and the derivative is positive.
\item Under the preparing rule, Lemma~\ref{lem:prices}(d) and \(\pi<\tau(r)\) give entry
\((\rho,\rho,\rho+(1-\rho)\chi_r)\) after the orders \((-1,0,1)\), because \(B_r(M)\ge c_H\) exactly when
\(M\ge\tau(r)\). By (iv) and (v), the class component \(\rho\DT(r)(M-m)\) and the sale impact \(\rho\DT(r)(\pi-m)\)
rise strictly with \(\DT\). The entry component is \((1-\rho)\chi_r\,w_r(M)\), the purchase impact is
\(\rho\DT(r)(M-\pi)+(1-\rho)\chi_r\,w_r(M)\), and by (ii) the weighted impact difference is
\((1-\rho)\chi_r\,\bar f(1)\,w_r(M)\). Where \(\chi_r=1\) each of these rises strictly, because \(w_r(M)\) does; where
\(\chi_r=0\) the entry component and the weighted impact difference are zero, and the class component, both impacts
and the price gap rise strictly with \(\DT\).
\item Each \(g_j\) is nonincreasing along the path (Lemma~\ref{lem:twonumbers}), so \(B_r(M)=(1-M)g_L(F_r)+Mg_H(F_r)\)
is nonincreasing, and \(\chi_r\) is nonincreasing. In \(g_j(F_r)=\int_p^\infty F_r(y)[1-G_j(y)]\,dy\) the integrand is
bounded by \(1-G_j\), which is integrable because \(G_j\) has a finite mean; when \(F_r(y)\) is continuous in \(r\),
dominated convergence makes \(B_r(M)\) continuous in \(r\), and \(\{r:B_r(M)\ge c_H\}\) is closed in the path. In the
same way, the display of (a) makes \(w_r(\mu)\) continuous in \(r\), and \(t_0=p[1-F_r(p)]\) is continuous in \(r\).
At \(r_C\), continuity and \(B_r(M)<c_H\) for \(r>r_C\) give \(B_{r_C}(M)=c_H\), that is \(M=\tau(r_C)\).

On the neighborhood of \(r_C\), \(P(-1)=t_0+\rho\,w_r(m)\) and \(P(0)=t_0+\rho\,w_r(\pi)\) are continuous in \(r\), and so
is \(I_-\). When an indifferent high-cost challenger prepares, \(e(1)=1\) for \(r\le r_C\) and \(e(1)=\rho\) for
\(r>r_C\), so \(P(1)=t_0+e(1)\,w_r(M)\) is continuous from the left at \(r_C\) and its right limit is lower by
\((1-\rho)\,w_{r_C}(M)>0\); \(I_+\) and the price gap move with \(P(1)\), and by (ii) the weighted impact difference
\(\bar f(1)\,[e(1)-\rho]\,w_r(M)\) falls by \((1-\rho)\,\bar f(1)\,w_{r_C}(M)\).

When an indifferent high-cost challenger may prepare with any probability, the posterior after a purchase at \(r_C\)
equals \(\tau(r_C)\), so (E3) allows every \(e(1)\in[\rho,1]\). The class \(L\) conditions and the posteriors do not
depend on \(e(1)\). On the right of \(r_C\), full two-sided trade has entry \(\rho\) after a purchase, so the purchase
margin \(\rho\DT(r)(1-M)-k\) is nonnegative there; its limit \(\rho\DT(r_C)(1-M)-k\) is nonnegative as well. At \(r_C\)
the purchase earns \(e(1)\DT(r_C)(1-M)-k\), which is positive for every \(e(1)\in(\rho,1]\) and nonnegative at
\(e(1)=\rho\), where an indifferent investor who may mix buys with probability one. The paragraph on constructed tuples
in Appendix~\ref{app:beliefs} completes each tuple. The price gap at \(r_C\) is
\(\rho\DT(r_C)(M-m)+[e(1)-\rho]\,w_{r_C}(M)\), affine in \(e(1)\), and it takes every value between its two
one-sided limits.\qedhere
\end{enumerate}
\end{enumerate}
\end{proof}

\begin{proposition}[What order counts and prices identify]\label{prop:informedshare}
Let full two-sided trade hold.
\begin{enumerate}[label=(\alph*),leftmargin=2em]
\item \emph{A known prior other than \(1/2\).} If \(\pi\ne1/2\) is known, the counts identify the slot:
\[
\lambda=\frac{\E[X]}{2\pi-1},\qquad b=\frac{v-\lambda}{2},\qquad \beta=\frac{v-\lambda}{2(1-\lambda)},\qquad
\vartheta=\frac{\lambda}{v}.
\]
Over independent order opportunities, the mean order divided by \(2\pi-1\) estimates \(\lambda\) without bias, with
variance \(\bigl(v-\E[X]^2\bigr)/(2\pi-1)^2\) divided by the number of opportunities. If also \(\pi<\tau\), the sale
impact identifies \(\rho\DT=\bar f(-1)\,I_-/[\pi(1-\pi)\lambda]\).
\item \emph{The prior \(1/2\).} At the prior \(1/2\) the order law is \((v/2,1-v,v/2)\), so the counts identify \(v\) and
no other function of the slot. A share \(\vartheta'\in(0,1)\) belongs to exactly one pair \((\lambda',\beta')\) with
the volume \(v\), namely \(\lambda'=v\vartheta'\) and \(\beta'=v(1-\vartheta')/[2(1-v\vartheta')]\), and this pair is
a slot of the model, \(\beta'\le1/3\), exactly when \(\vartheta'\ge3-2/v\). Let \(\tau>1/2\). Under full two-sided
trade at this slot, entry is \(\rho+(1-\rho)(v/2)\,\1\{(1+\vartheta')/2\ge\tau\}\) when an indifferent high-cost
challenger prepares. Full two-sided trade at this slot is an equilibrium exactly when
\(\vartheta'\le1-2k/(\rho\DT)\) if an indifferent investor may mix over its best orders, and exactly when
\(\vartheta'<1-2k/(\rho\DT)\) if an indifferent investor abstains.
\item \emph{The sale price.} If entry after a sale and after no order is \(\rho\) and \(\rho\DT\) is known, the sale
impact gives \(\pi(1-\pi)\lambda=\bar f(-1)\,I_-/(\rho\DT)\), and with the counts
\[
\lambda=\frac{2\bar f(-1)\,I_-}{\rho\DT}+\sqrt{\Bigl(\frac{2\bar f(-1)\,I_-}{\rho\DT}\Bigr)^2+\E[X]^2},\qquad
\pi=\frac12\Bigl(1+\frac{\E[X]}{\lambda}\Bigr),
\]
and then \(b\) and \(\beta\) as in part (a). At the prior \(1/2\), \(\lambda=4\bar f(-1)\,I_-/(\rho\DT)\) and
\(\vartheta=2I_-/(\rho\DT)\).
\end{enumerate}
\end{proposition}

\begin{proof}[Proof of Proposition~\ref{prop:informedshare}]
Let full two-sided trade hold. The counts are the order law \(\bar f\).
\begin{enumerate}[label=(\alph*),leftmargin=2em]
\item By \eqref{eq:a8-full}, \(v=2b+\lambda\) and \(\E[X]=\lambda(2\pi-1)\). With \(\pi\ne1/2\) known, these give
\(\lambda\), then \(b=(v-\lambda)/2\), \(\beta=b/(1-\lambda)\) and \(\vartheta=\lambda/v\). The order \(X\) takes values in
\(\{-1,0,1\}\), so \(X^2=\1\{X\ne0\}\), \(\E[X^2]=v\) and the variance of \(X\) is \(v-\E[X]^2\). The mean of a sample of
independent orders, divided by \(2\pi-1\), has mean \(\lambda\) and variance \((v-\E[X]^2)/(2\pi-1)^2\) divided by the
sample size. If also
\(\pi<\tau\), entry after a sale and after no order is \(\rho\), so \(I_-=\rho\,[w_F(\pi)-w_F(m)]=\rho\DT\,(\pi-m)\), and
\eqref{eq:a8-balance} gives \(I_-=\rho\DT\,\pi(1-\pi)\lambda/\bar f(-1)\). Solving for \(\rho\DT\) gives the formula.
\item At the prior \(1/2\), \eqref{eq:a8-full} gives \(\bar f(\pm1)=b+\lambda/2=v/2\) and \(\bar f(0)=1-v\), so the order
law is a function of \(v\) alone. Given \(v\) and \(\vartheta'\), the equations \(\lambda'=v\vartheta'\) and
\(2b'+\lambda'=v\) give \(b'=v(1-\vartheta')/2\) and \(\beta'=b'/(1-\lambda')=v(1-\vartheta')/[2(1-v\vartheta')]\), a unique
pair. Since \(v=1-\bar f(0)<1\), \(0<\lambda'<1\) and \(\beta'>0\). The inequality \(\beta'\le1/3\) is
\(3v(1-\vartheta')\le2(1-v\vartheta')\), that is \(v\vartheta'\ge3v-2\), that is \(\vartheta'\ge3-2/v\).

At this slot and the prior \(1/2\), \(M'=(1+\vartheta')/2\) and \(m'=(1-\vartheta')/2\) by the standing paragraph. Let
\(\tau>1/2\). Under full two-sided trade, the posteriors after a sale and after no order are \(m'<1/2<\tau\), so
entry after them is \(\rho\), and when an indifferent high-cost challenger prepares, entry after a purchase is one
exactly when \(M'\ge\tau\) and \(\rho\) otherwise. Weighting by \(\bar f(1)=v/2\) gives the entry formula. By
Lemma~\ref{lem:trading}, the sale earns \(K_L-k=\rho\DT\,m'-k\), and the purchase earns
\(K_H-k=e(1)\DT\,(1-M')-k\ge\rho\DT\,m'-k\), because \(1-M'=m'\) and \(e(1)\ge\rho\). Wrong-signed orders never pay. When
an indifferent investor may mix over its best orders, full orders are optimal exactly when \(K_L\ge k\) and
\(K_H\ge k\), that is exactly when \(\rho\DT(1-\vartheta')/2\ge k\); when an indifferent investor abstains, the two
inequalities are strict. The first condition is \(\vartheta'\le1-2k/(\rho\DT)\), and the second the same with strict
inequality. The paragraph on constructed tuples in Appendix~\ref{app:beliefs} completes the tuple under either
challenger rule, so these are also the conditions for an equilibrium.
\item With entry \(\rho\) after a sale and after no order, \(I_-=\rho\DT\,(\pi-m)\), and \eqref{eq:a8-balance} gives
\(q:=\bar f(-1)\,I_-/(\rho\DT)=\pi(1-\pi)\lambda\). Since \((2\pi-1)^2=1-4\pi(1-\pi)\),
\[
\E[X]^2=\lambda^2(2\pi-1)^2=\lambda^2-4q\lambda,\qquad\text{so}\qquad \lambda^2-4q\lambda-\E[X]^2=0 .
\]
Since \(q>0\), the product of the two roots is \(-\E[X]^2\le0\) and their sum \(4q\) is positive, so exactly one root is
positive, \(\lambda=2q+\sqrt{4q^2+\E[X]^2}\), which is the display. Then \(2\pi-1=\E[X]/\lambda\), and part (a)
gives \(b\) and \(\beta\) from \(v\) and \(\lambda\). At the prior \(1/2\), \(\E[X]=0\), so \(\lambda=4q\), and
\(\vartheta=\lambda/v=2I_-/(\rho\DT)\) because \(\bar f(-1)=v/2\).\qedhere
\end{enumerate}
\end{proof}

\begin{corollary}[Counts under purchases only]\label{cor:purchasesonly}
Take an equilibrium in which class \(L\) abstains and class \(H\) buys with probability \(\sigma_H(1)>0\), and write
\(s_+=\lambda\sigma_H(1)\).
\begin{enumerate}[label=(\alph*),leftmargin=2em]
\item Then \(\bar f(-1)=b\), \(\bar f(1)=b+\pi s_+\), \(\E[X]=\pi s_+\) and \(v=2b+\pi s_+\). At every prior, the
counts identify \(b=(v-\E[X])/2\) and the share \(\E[X]/v\) of nonzero orders that the informed investor places. At
a known prior they identify \(s_+=\E[X]/\pi\); if class \(H\) buys with probability one, \(\lambda=\E[X]/\pi\).
\item Suppose that an indifferent investor may mix over its best orders and that class \(H\) mixes,
\(0<\sigma_H(1)<1\). For every \(\lambda'\in[s_+,1-3b]\), the slot \((\lambda',\beta')\) with \(\beta'=b/(1-\lambda')\), with
\(\sigma'_H(1)=s_+/\lambda'\) and with class \(L\) abstaining, has an equilibrium with the same class-conditional
order laws, and so the same posteriors, prices, entry and acquisition probabilities. The counts
therefore identify \(b\) and, at a known prior, \(s_+\); even with the prices, they do not identify \(\lambda\). On a mixed purchase, the indifference of
class \(H\) gives \(s_+=(b/\pi)\bigl[e(1)(1-\pi)\DT/k-1\bigr]\), which is \((b/\pi)\bigl[(1-\pi)\DT/k-1\bigr]\) when
every challenger prepares after a purchase.
\end{enumerate}
\end{corollary}

\begin{proof}[Proof of Corollary~\ref{cor:purchasesonly}]
By Lemma~\ref{lem:trading}(a) no class places a wrong-signed order, so class \(H\) places the order \(0\) with
probability \(1-\sigma_H(1)\). The class-conditional order laws on \((-1,0,1)\) are
\[
f_H=(b,\,1-2b-s_+,\,b+s_+),\qquad f_L=(b,\,1-2b,\,b).
\]
\begin{enumerate}[label=(\alph*),leftmargin=2em]
\item Averaging with the weights \(\pi\) and \(1-\pi\) gives \(\bar f(-1)=b\) and \(\bar f(1)=b+\pi s_+\), so
\(\E[X]=\pi s_+\) and \(v=2b+\pi s_+\). Hence \(b=(v-\E[X])/2\). The informed investor places a nonzero order with
probability \(\pi s_+=\E[X]\), so its share of nonzero orders is \(\E[X]/v\). At a known prior, \(s_+=\E[X]/\pi\), and
\(s_+=\lambda\) when \(\sigma_H(1)=1\).
\item The original slot has \(\beta\le1/3\), that is \(\lambda\le1-3b\), and \(s_+=\lambda\sigma_H(1)<\lambda\), so the
segment \([s_+,1-3b]\) is nondegenerate and contains \(\lambda\). Take \(\lambda'\) in it. Then \(\beta'=b/(1-\lambda')\)
lies in \((0,1/3]\) because \(1-\lambda'\ge3b\), the liquidity mass is \((1-\lambda')\beta'=b\), and
\(\sigma'_H(1)=s_+/\lambda'\) lies in \((0,1]\) because \(\lambda'\ge s_+\). The class-conditional order laws of this slot
are \(f_H\) and \(f_L\) above, since they depend only on \(b\) and \(s_+\). By \eqref{eq:odds} the posteriors are the
same, so keep the original entry \(e(x)\) after each order: it meets Lemma~\ref{lem:prices}(d) at these posteriors
under either challenger rule. Prices \(t_0+e(x)\,w_F(\mu(x))\), the preparation rule and the price beliefs are then the
same as well. The payoffs of Lemma~\ref{lem:trading} depend only on posteriors and entry, so they do not change, and
by Lemma~\ref{lem:trading}(d) the factor \(\lambda'\) changes no comparison. In the original equilibrium class \(H\)
mixes, so the purchase and abstention are its two best orders, and an indifferent investor who may mix over them can
buy with probability \(\sigma'_H(1)\); abstention stays optimal for class \(L\). So the slot \((\lambda',\beta')\) has an
equilibrium with the same order laws, posteriors, prices and entry, and by Proposition~\ref{prop:winner}(i) the same
acquisition probability \(e(x)\,W_F(\mu(x))\) after each order. The counts and the prices are therefore the same at
every \(\lambda'\) in the segment; part (a) gives \(b\) and, at a known prior, \(s_+\).

On a mixed purchase, class \(H\) is indifferent, \(K_H=e(1)\DT\,[1-\mu(1)]=k\), with
\(1-\mu(1)=(1-\pi)b/(b+\pi s_+)\) by Bayes' rule. Solving \(e(1)\DT(1-\pi)b=k\,(b+\pi s_+)\) for \(s_+\) gives the
formula, and \(e(1)=1\) gives its second form.\qedhere
\end{enumerate}
\end{proof}

Let \(\1_{\rm prep}\) and \(\1_{\rm acq}\) indicate that the challenger prepares and that it acquires the target. For
a class \(j\), \(W_j(F)=\Pr(R<\theta)\) with \(\theta\) drawn from \(G_j\) is the probability that a prepared
class-\(j\) challenger outbids the incumbent. Write \(\Omega(F)=W_H(F)-W_L(F)\) and \(W_F(\mu)=W_L(F)+\mu\,\Omega(F)\).

\begin{proposition}[Orders predict the winner]\label{prop:winner}
\leavevmode
\begin{enumerate}[label=(\roman*),leftmargin=2em]
\item In every equilibrium, the probability that the challenger acquires the target after the order \(x\), given
preparation, is \(W_F(\mu(x))\), and \(\E[\1_{\rm acq}\mid X=x]=e(x)\,W_F(\mu(x))\). Also
\(\Omega(F)=\E_F[D(R)]\ge0\), and \(\Omega(F)>0\) exactly when \(\Pr_F\{D(R)>0\}>0\).
\item Under full two-sided trade, \(W_F(M)-W_F(m)=(M-m)\,\Omega(F)\ge0\), positive exactly when
\(\Pr_F\{D(R)>0\}>0\). The prices rise strictly with the order, so the win probability of a prepared challenger
rises with the price, strictly when \(\Pr_F\{D(R)>0\}>0\).
\item Under full two-sided trade with \(\pi<\tau\),
\[
\begin{aligned}
\Cov(X,\1_{\rm prep})&=[e(1)-\rho]\,\bar f(1)\,(1-\E[X]),\\
\Cov(X,\1_{\rm acq})&=2\pi(1-\pi)\rho\lambda\,\Omega(F)+[e(1)-\rho]\,\bar f(1)\,(1-\E[X])\,W_F(M).
\end{aligned}
\]
If \(B_F(M)<c_H\), entry is \(\rho\) after every order, \(\Cov(X,\1_{\rm prep})=0\), and
\(\Cov(X,\1_{\rm acq})=2\pi(1-\pi)\rho\lambda\,\Omega(F)\), which is positive exactly when \(\Pr_F\{D(R)>0\}>0\).
\item Under full two-sided trade with \(\pi<\tau\),
\[
\begin{aligned}
&\E[\1_{\rm acq}\mid X=1]-2\,\E[\1_{\rm acq}\mid X=0]+\E[\1_{\rm acq}\mid X=-1]\\
&\qquad=\rho\,(M+m-2\pi)\,\Omega(F)+[e(1)-\rho]\,W_F(M).
\end{aligned}
\]
At the prior \(1/2\), \(M+m=1=2\pi\). If also an indifferent high-cost challenger prepares, then
\(\Cov(X,\1_{\rm acq})=\rho\lambda\,\Omega(F)/2+(1-\rho)\,\chi_F\,\bar f(1)\,W_F(M)\), and the second difference
equals \((1-\rho)\,\chi_F\,W_F(M)\), which is positive exactly when a purchase recruits.
\item Let \(F_r\) be a path, hold the prior and the slot fixed, and let class \(H\) always outbid the incumbent:
\(W_H(F_r)=1\) at every strength. At strengths with full two-sided trade, the winner gap
\(W_{F_r}(M)-W_{F_r}(m)=(M-m)\,[1-W_L(F_r)]\) is nondecreasing in \(r\), and it is larger at \(r'\) than at
\(r<r'\) exactly when \(W_L(F_{r'})<W_L(F_r)\).
\item If \(F_r\) is uniform on \([0,r]\) with \(r\ge p\), then \(\Omega(F_r)=\Gamma(r)/r\), and its right derivative
in \(r\) is \([D(r)-\Omega(F_r)]/r\). So \(\Omega(F_r)\), and with it the full-trade winner gap
\((M-m)\,\Omega(F_r)\) at a fixed prior and slot, is nondecreasing on an interval of strengths exactly when
\(D(r)\ge\Omega(F_r)\) at almost every strength of the interval, and it rises strictly on an interval where
\(D(r)>\Omega(F_r)\) at almost every strength.
\end{enumerate}
\end{proposition}

\begin{proof}[Proof of Proposition~\ref{prop:winner}]
\begin{enumerate}[label=(\roman*),leftmargin=2em]
\item By Step 0 of the proof of Lemma~\ref{lem:prices}, given the class \(j\) and the order \(x\) the challenger
prepares with probability \(e(x)\), independently of \((\theta,R)\). Preparation therefore leaves the class posterior
at \(\mu(x)\), and given the class it leaves \((\theta,R)\) at its prior law. Bids are truthful by (E6), and every value
meets the reserve, so a prepared challenger acquires the target exactly when \(R<\theta\); ties have probability zero
because \(F\) is continuous. Given the class \(j\) and preparation, it acquires with probability \(W_j(F)\), and given
preparation alone with probability \((1-\mu(x))W_L(F)+\mu(x)W_H(F)=W_F(\mu(x))\). Without preparation it does not
acquire, so \(\E[\1_{\rm acq}\mid X=x]=e(x)\,W_F(\mu(x))\). Also \(W_j(F)=\int\Pr_j(\theta>y)\,dF(y)=\int[1-G_j(y)]\,dF(y)\),
since \(\Pr_j(\theta>y)\) and \(1-G_j(y)\) differ at countably many \(y\), which \(F\) does not charge. Hence
\(\Omega(F)=\int[G_L(y)-G_H(y)]\,dF(y)=\E_F[D(R)]\). The class ranking gives \(D\ge0\), so \(\Omega(F)\ge0\), with
\(\Omega(F)>0\) exactly when \(D(R)>0\) with positive probability.
\item Under full two-sided trade the posteriors after a purchase and a sale are \(M\) and \(m\), and \(W_F\) is affine
with slope \(\Omega(F)\), so \(W_F(M)-W_F(m)=(M-m)\,\Omega(F)\), with \(M>m\). By Lemma~\ref{lem:prices}(e) the prices
rise strictly with the posterior, so \(P(-1)<P(0)<P(1)\), and the win probability \(W_F(\mu(x))\) of a prepared
challenger is a nondecreasing function of the price, strictly increasing when \(\Omega(F)>0\).
\item For any indicator \(Y\), \(\Cov(X,Y)=\Cov(X,\E[Y\mid X])\). With \(\pi<\tau\), entry after the orders
\((-1,0,1)\) is \((\rho,\rho,e(1))\), so \(\E[\1_{\rm prep}\mid X]=\rho+[e(1)-\rho]\,\1\{X=1\}\). For any number \(c\),
\(\Cov(X,c\,\1\{X=1\})=c\,[\bar f(1)-\E[X]\,\bar f(1)]=c\,\bar f(1)\,(1-\E[X])\), which gives the first line. By (i),
\(\E[\1_{\rm acq}\mid X]=\rho\,W_F(\mu(X))+[e(1)-\rho]\,W_F(M)\,\1\{X=1\}\), and
\(\Cov(X,\rho\,W_F(\mu(X)))=\rho\,\Omega(F)\Cov(X,\mu(X))\). Since \(\mu(X)=\Pr(H\mid X)\),
\(\E[X\mu(X)]=\E[X\,\1\{j=H\}]=\pi\,\E[X\mid H]=\pi\lambda\) under full two-sided trade, so
\(\Cov(X,\mu(X))=\pi\lambda-\lambda(2\pi-1)\pi=2\pi(1-\pi)\lambda\). Adding the two terms gives the second line. If
\(B_F(M)<c_H\), then \(M<\tau\), every posterior lies below \(\tau\), entry is \(\rho\) after every order, and
\(e(1)=\rho\) in both lines; positivity follows from (i).
\item By (i), with entry \((\rho,\rho,e(1))\) and posteriors \((m,\pi,M)\), the second difference is
\(e(1)W_F(M)-2\rho W_F(\pi)+\rho W_F(m)=\rho\,[W_F(M)-2W_F(\pi)+W_F(m)]+[e(1)-\rho]\,W_F(M)\), and
\(W_F(M)-2W_F(\pi)+W_F(m)=(M+m-2\pi)\,\Omega(F)\) because \(W_F\) is affine. At the prior \(1/2\), \(M+m=1\) by the
standing paragraph and \(\E[X]=0\). When an indifferent high-cost challenger prepares, \(e(1)=\rho+(1-\rho)\chi_F\),
because \(B_F(M)\ge c_H\) exactly when \(M\ge\tau\). Substituting in (iii) and in the second difference gives the two
expressions. If \(\chi_F=1\), then \((1-M)g_L+Mg_H=B_F(M)\ge c_H>0\), so some class has positive profit
\(\E_j[(\theta-\max\{p,R\})_+]\), which is positive only when \(\Pr(R<\theta)>0\); that class has \(W_j(F)>0\), and
\(W_F(M)\), which puts positive weight on both classes, is positive. So the second difference is positive exactly when
\(\chi_F=1\), that is, exactly when a purchase recruits.
\item With \(W_H(F_r)=1\), \(\Omega(F_r)=1-W_L(F_r)\), and (ii) gives the winner gap. Since
\(W_L(F_r)=\E_L[F_r(\theta)]\) and \(F_{r'}\le F_r\) for \(r'>r\), \(W_L\) is nonincreasing along the path. The prior and
the slot are fixed, so \(M-m>0\) is fixed, and the gap at \(r'\) exceeds the gap at \(r\) exactly when
\(W_L(F_{r'})<W_L(F_r)\).
\item Write \(\Gamma(x)=\int_p^xD(y)\,dy\) for \(x\ge p\), the class difference in expected proceeds at the incumbent
value \(x\) (Appendix~\ref{app:payoffs}). Under the class ranking \(G_H=G_L=0\) below \(p\), so \(D=0\) there, and for
the uniform law on \([0,r]\) with \(r\ge p\), \(\Omega(F_r)=\E_{F_r}[D(R)]=\frac1r\int_0^rD(y)\,dy=\Gamma(r)/r\). Since
\(D\) is right-continuous, \(\Gamma\) has the right derivative \(D(r)\), and the quotient rule gives the right
derivative \([rD(r)-\Gamma(r)]/r^2=[D(r)-\Omega(F_r)]/r\). Since \(0\le D\le1\), \(\Gamma\) is Lipschitz, so
\(\Omega(F_r)\) is absolutely continuous on every compact interval of positive strengths, and its derivative equals
\([D(r)-\Omega(F_r)]/r\) at every continuity point of \(D\), hence almost everywhere, because \(D\) is a difference of
two monotone functions. An absolutely continuous function is nondecreasing on an interval exactly when its derivative is
nonnegative almost everywhere there, and it rises strictly where its derivative is positive almost everywhere. The
prior and the slot are fixed, so the full-trade winner gap \((M-m)\,\Omega(F_r)\) inherits both properties.\qedhere
\end{enumerate}
\end{proof}

Add a public event after the price that reveals only whether the challenger prepared, and nothing else. Market makers
then price again.

\begin{proposition}[Announcement returns]\label{prop:announcement}
Take any equilibrium, at any prior, under either challenger rule.
\begin{enumerate}[label=(\roman*),leftmargin=2em]
\item After the order \(x\), the price changes by \(\Delta^1_x=[1-e(x)]\,w_F(\mu(x))\) when the event reports
preparation, and by \(\Delta^0_x=-e(x)\,w_F(\mu(x))\) when it reports no preparation. The second report has positive
probability exactly when \(e(x)<1\).
\item After an order \(x\) with \(e(x)<1\), the mean change over both reports, with their probabilities, is zero.
\item Pool all orders, and take the sample of preparation reports. Its mean price change is
\[
\sum_x\Pr(X=x\mid\text{preparation})\,\Delta^1_x
=\frac{\sum_x\bar f(x)\,e(x)\,[1-e(x)]\,w_F(\mu(x))}{\sum_x\bar f(x)\,e(x)},
\]
which is positive exactly when entry is below one after some order, that is, exactly when preparation is uncertain.
If \(B_F(\pi)<c_H\), entry is \(\rho\) after some order, and the mean is positive.
\item The expected price before the order opportunity minus the price after no order is the weighted impact
difference, \(\E[P(X)]-P(0)=\bar f(1)I_+-\bar f(-1)I_-\), which parts (i) to (iii) of Proposition~\ref{prop:spread}
evaluate. Under full two-sided trade the posterior after no order is the prior.
\end{enumerate}
\end{proposition}

\begin{proof}[Proof of Proposition~\ref{prop:announcement}]
After the order \(x\), the event reveals whether the challenger prepared, and the market makers price at the expected
value of the target given the order and the report.
\begin{enumerate}[label=(\roman*),leftmargin=2em]
\item By the proof of Proposition~\ref{prop:winner}(i), preparation leaves the class posterior at \(\mu(x)\) and, given
the class \(j\), leaves \((\theta,R)\) at its prior law. With truthful bids (E6), expected proceeds in class \(j\) after
preparation are \(t_j\), so the price after a report of preparation is \((1-\mu(x))t_L+\mu(x)t_H=t_0+w_F(\mu(x))\).
Without preparation the target receives the proceeds without entry, with mean \(t_0\) in both classes. Subtracting
\(P(x)=t_0+e(x)\,w_F(\mu(x))\) gives \(\Delta^1_x\) and \(\Delta^0_x\). Given the order \(x\), the report of no
preparation has probability \(1-e(x)\), which is positive exactly when \(e(x)<1\).
\item When \(e(x)<1\), \(e(x)\Delta^1_x+[1-e(x)]\Delta^0_x=e(x)[1-e(x)]\,w_F(\mu(x))-[1-e(x)]\,e(x)\,w_F(\mu(x))=0\).
\item The probability of the order \(x\) given preparation is \(\bar f(x)e(x)/\sum_{x'}\bar f(x')e(x')\), which gives the
display. Every order has \(\bar f(x)>0\) by Lemma~\ref{lem:beliefs}(a), entry \(e(x)\ge\rho>0\), and an interior
posterior, so \(w_F(\mu(x))>0\). Every term of the numerator is therefore nonnegative, and the numerator is positive
exactly when \(e(x)<1\) after some order. Every order has positive probability, so this is exactly when preparation
is uncertain. If \(B_F(\pi)<c_H\), Bayes' rule and the affine form of \(B_F\) give
\(\sum_x\bar f(x)\,B_F(\mu(x))=B_F(\pi)<c_H\), so some order has \(B_F(\mu(x))<c_H\), that is \(\mu(x)<\tau\), and
Lemma~\ref{lem:prices}(d) gives \(e(x)=\rho<1\) after it, under both challenger rules.
\item The identity is part (i) of Proposition~\ref{prop:spread}, which holds in every equilibrium. Under full
two-sided trade, Lemma~\ref{lem:beliefs}(d) gives \(\mu(0)=\pi\).\qedhere
\end{enumerate}
\end{proof}

\begin{named}{The example}
Take the example at the prior \(1/2\): the incumbent is uniform on \([0,r]\), the classes are point masses at \(h\)
and \(\ell\), and the slot lies in the band \(\tau(r_1)<M<\tau(r_2)\) of Corollary~\ref{prop:region}. Unless a part
names another tie rule, an indifferent investor may mix over its best orders and an indifferent high-cost challenger
prepares. The numbers are those of the example's slot.
\begin{enumerate}[label=(\alph*),leftmargin=2em]
\item \emph{Full two-sided trade.} For \(1<r<h\), \(B_r(1/2)=5-r/4+1/(8r)<c_H\), so \(\tau(r)>1/2\). The strength
\(r_S\) solves \(\rho\DT(r)\,m=k\): \(r_S=\fr\bigl(k/(\rho m)\bigr)\) with \(k/(\rho m)=\reg{delta_rS_B}\), and
\(r_S\approx\reg{rS_B}\). For \(r_S<r<h\), every equilibrium has full two-sided trade, under both investor rules and
both challenger rules, whether or not the challenger sees the price. Its posteriors are \((m,1/2,M)\); its order
frequencies are \(\bar f(\pm1)=b+\lambda/2=\reg{d_slot_B}\) and \(\bar f(0)=1-2b-\lambda\); its volume is
\(v=\reg{vol_full_B}\); \(\E[X]=0\); and the informed share is \(\vartheta=M-m=\reg{theta_B}\). By
Lemma~\ref{lem:a4-order}, \(r_S<r_1<r_C<r_2<h\), and \(r_C-r_S=\reg{margin_obs_rS_rC_B}\).
\item \emph{Recruitment and entry.} The top profit \(B_r(M)\) falls strictly in \(r\) for \(r>\ell\), equals
\(9M+3/8\) at \(r=\ell\), and lies below \(c_H\) at \(r=h\). In the band, \(M>\tau(r_1)=\reg{tau_r1_frac}>5/8\), so
the entry ceiling \(r_C\) is the unique strength with \(B_r(M)=c_H\), and \(\chi_r=1\) exactly when \(r\le r_C\).
When an indifferent high-cost challenger prepares, entry after the orders \((-1,0,1)\) is
\((\rho,\rho,\rho+(1-\rho)\chi_r)\) on \((r_S,h)\), and entry is \(\rho+(1-\rho)\chi_r\,(b+\lambda/2)\). From \(r_0\),
where the orders are zero and entry is \(\rho\), to \(r_1\), entry rises by \((1-\rho)(b+\lambda/2)=\reg{entry_jump_B}\).
With a hidden price, entry is \(\rho\) after every order.
\item \emph{The price gap.} Here \(w_r(\mu)=1/2-1/(8r)+\DT(r)\,\mu\) with \(\DT(r)=(r-1)^2/(2r)\), and both
\(1/2-1/(8r)\) and \(\DT(r)\) rise strictly in \(r\) for \(r>\ell\); also, for \(r\ge\ell\),
\(\min\{\ell,r\}^2-3p^2=\ell^2-3p^2=\reg{margin_reserve_unif_B}>0\), the condition of
Proposition~\ref{prop:spread}(vi)(b). At the prior \(1/2\), \(\bar f(1)=\bar f(-1)\), and the weighted test reads
\(I_+-I_-=(1-\rho)\chi_r\,w_r(M)\), so \(I_+>I_-\) exactly when \(r\le r_C\). On \((r_S,r_C]\), both components of
the price gap, both impacts, the price gap and the weighted impact difference rise strictly with \(r\). On
\((r_C,h)\), the entry component and the weighted impact difference are zero, and the class component, both impacts
and the price gap rise strictly. At \(r_1\) the class component is \(\reg{ba_class_r1_B}\) and the entry component
\(\reg{ba_entry_r1_B}\), so the price gap is \(\reg{ba_r1_B}\), of which the entry component is
\(\reg{ba_share_r1_B}\); the weighted impact difference \(\E[P(X)]-P(0)\) is \(\reg{wtest_r1_B}\). With a hidden price,
the price gap at \(r_1\) is the class component \(\reg{ba_class_r1_B}\). At \(r_C\) the price gap falls from
\(\reg{ba_rC_left_B}\) to \(\reg{ba_rC_right_B}\), by \(\reg{ba_jump_rC_B}\), and at \(r_2\) it is the class component
\(\reg{ba_r2_B}\). When an indifferent investor may mix over its best orders and an indifferent high-cost challenger
may prepare with any probability, entry after a purchase
at \(r_C\) can take any value in \([\rho,1]\), and the price gap at \(r_C\) any value from \(\reg{ba_rC_right_B}\)
to \(\reg{ba_rC_left_B}\).
\item \emph{Identification.} The counts identify only the volume. When an indifferent investor may mix over its best
orders, the informed shares that the counts and full two-sided trade at \(r_1\) allow form the interval from
\(3-2/v\approx\reg{theta_min_counts_B}\) to \(1-2k/(\rho\DT(r_1))=\reg{theta_max_full_eq_r1_B}\), both ends
included; when an indifferent investor abstains, the right end is excluded. At \(r_2\) the right end is
\(\reg{theta_max_full_eq_r2_B}\). The shares with strict recruitment at \(r_1\) and strict exclusion at \(r_2\) form
the band \((2\tau(r_1)-1,\,2\tau(r_2)-1)=(\reg{theta_lo},\reg{theta_hi})\). Each share in the band belongs to a slot
with the example's counts and the same entry at \(r_1\) and at \(r_2\), and full two-sided trade is the unique outcome
of that slot at both strengths. The slot with share \(\reg{alt_theta_B}\) is one of them; its price gap at \(r_1\)
is \(\reg{alt_ba_r1_B}\), against \(\reg{ba_r1_B}\). On \((r_S,h)\) the sale impact identifies the share,
\(\vartheta=2I_-/(\rho\DT(r))\).
\item \emph{Purchases only.} A full purchase with class \(L\) abstaining is an equilibrium exactly on \([r_F,r_N]\),
and a mixed purchase that recruits is an equilibrium on \((r_A,r_F)\), each beside a quiet market
(Corollary~\ref{cor:example-coexistence}). On \([r_F,r_N]\) the counts identify \(\lambda=2\E[X]\). On \((r_A,r_F)\),
\(s_+=b\,[\DT(r)/k-2]\), and the counts and the prices do not identify \(\lambda\).
\item \emph{Winners.} For \(\ell\le r\le h\), \(W_H=1\), \(W_L=\ell/r\) and \(\Omega=1-\ell/r\), so a prepared
challenger wins after the order \(x\) with probability \(\mu(x)+[1-\mu(x)]\,\ell/r\). Under full two-sided trade the
winner gap \(\vartheta(1-\ell/r)\) rises strictly with \(r\): class \(H\) always outbids the incumbent and \(W_L\)
falls strictly, and \(D(r)-\Omega=\ell/r>0\) for \(\ell\le r<h\), which is \(\reg{margin_winsep_r1_B}\) at \(r_1\) and
\(\reg{margin_winsep_r2_B}\) at \(r_2\). The winner gap is \(\reg{wgap_r1_B}\) at \(r_1\) and \(\reg{wgap_r2_B}\) at
\(r_2\). The covariance of the order with entry is \(\reg{cov_entry_r1_B}\) at \(r_1\) and zero at \(r_2\). Its
covariance with a challenger win is \(\reg{cov_win_r1_B}\) at \(r_1\), the class term
\(\rho\lambda\Omega/2=\reg{cov_win_class_r1_B}\) plus the recruitment term \(\reg{cov_win_recruit_r1_B}\), and
\(\reg{cov_win_r2_B}\) at \(r_2\), the class term alone. The second difference of the win probability across the
orders is \(\reg{wsecond_r1_B}\) at \(r_1\) and zero at \(r_2\).
\item \emph{Announcements.} Since \(B_r(1/2)\) falls in \(r\), \(c_H-B_r(1/2)\ge c_H-B_{r_0}(1/2)=\reg{margin_res_half_B}\)
for every \(r\ge r_0\), so a sample of preparation reports has a positive mean price change in every equilibrium at
every such strength. At \(r_1\) the mean is \(\reg{evprep_mean_r1_B}\); the expected price before the order
opportunity is \(\reg{pbar_r1_B}\), and the price after no order is \(\reg{p_neutral_r1_B}\), lower by
\(\reg{wtest_r1_B}\). At \(r_2\) the mean is \(\reg{evprep_mean_r2_B}\), and the price after no order equals the
expected price.
\item \emph{Two clusters.} Let \(r_F<r<r_N\), so \(r<r_C\). When an indifferent investor may mix over its best orders
and an indifferent high-cost challenger prepares, Corollary~\ref{cor:example-coexistence} gives exactly two equilibria: a
quiet market, and a full purchase with class \(L\) abstaining. In the second, \(\bar f(1)=b+\lambda/2\) and
\(\bar f(-1)=b\), so \(\E[X]=\lambda/2=\reg{imb_buy_B}\). The posterior after a sale is \(1/2\), and after no order it
is \(\mu(0)=(1-2b-\lambda)/(2-4b-\lambda)\), with \(1/2-\mu(0)=\lambda/[4(1-v)]>0\) for the volume \(v=2b+\lambda/2\)
of this equilibrium. Both orders have entry \(\rho\), so \(P(-1)-P(0)=\rho\DT\lambda/[4(1-v)]>0\), which is
\(\reg{U_gap_r160_B}\) at \(r=\reg{str_r160}\); after a purchase every challenger prepares. In the quiet market the
three prices coincide and entry after a purchase is \(\rho\). When an indifferent high-cost challenger may prepare
with any probability, the threshold pair adds a third equilibrium: at \(r=\reg{str_r160}\), class \(H\) buys with
probability \(\reg{mix_sigma_r160_B}\), entry after a purchase is \(\reg{mix_e1_r160_B}\), the mean order is
\(\reg{mix_mean_order_r160_B}\), the low-class sale earns \(\reg{mix_lowsale_r160_B}\), and entry is
\(\reg{mix_entry_r160_B}\).
\end{enumerate}
\end{named}

\begin{proof}[Proof of the statements on the example]
At the prior \(1/2\), \(1-M=m\), so \(\kappa(1/2)=m\), and \(\bar f(\pm1)=b+\lambda/2\) under full two-sided trade.
For \(p<\ell<r<h\), Appendix~\ref{app:payoffs} gives
\[
t_0=p\Bigl(1-\frac pr\Bigr),\quad t_L=\ell-\frac{\ell^2-p^2}{2r},\quad
g_H=h-\frac r2-\frac{p^2}{2r},\quad g_L=\frac{\ell^2-p^2}{2r},\quad \DT(r)=\frac{(r-\ell)^2}{2r},
\]
and \(\DT\) is zero at \(\ell\) and strictly increasing above \(\ell\).
\begin{enumerate}[label=(\alph*),leftmargin=2em]
\item Gross profit at the prior is \(B_r(1/2)=(g_H+g_L)/2=h/2-r/4+(\ell^2-2p^2)/(4r)\), which is \(5-r/4+1/(8r)\) at the
example's values. It falls strictly in \(r\), and at \(r=\ell\) it equals \(h/2-p^2/(2\ell)<h/2<c_H\), so
\(B_r(1/2)<c_H\) for \(\ell<r<h\). There \(\gamma=g_H-g_L=h-r/2-\ell^2/(2r)>h-r>0\), so \(\tau(r)>1/2\). The strength
\(r_S\) is the unique strength above \(\ell\) at which \(\rho\DT(r)\,m=k\), and for \(r>r_S\),
\(\rho\kappa(1/2)\DT(r)=\rho m\DT(r)>k\). Proposition~\ref{prop:spread}(iv) gives full two-sided trade in every
equilibrium for \(r_S<r<h\), under both investor rules, whether or not the challenger sees the price; the challenger's
rule enters no step. The standing paragraph at the prior \(1/2\) gives the posteriors, the frequencies, the volume, the
mean order and the share. Lemma~\ref{lem:a4-order} gives \(r_S<r_1<r_C<r_2\), and \(r_2<h\).
\item By Appendix~\ref{app:payoffs}, \(\partial B_r(\mu)/\partial r<0\) for every \(\mu\) when \(\ell<r<h\), so
\(B_r(M)\) falls strictly there. At \(r=\ell\), \(g_L=(\ell^2-p^2)/(2\ell)\) and \(g_H=h-\ell/2-p^2/(2\ell)\), which give
\(9M+3/8\) at the example's values; it exceeds \(c_H\) exactly when \(M>5/8\). At \(r=h\),
\(B_h(M)<g_H=h/2-p^2/(2h)<h/2<c_H\), because \(g_L<g_H\) and \(M<1\). In the band, \(M>\tau(r_1)\), and
\(\tau(r_1)=\reg{tau_r1_frac}>5/8\). So \(B_r(M)\) crosses \(c_H\) once on \((\ell,h)\), at \(r_C\). Since \(B_r(M)\) is
nonincreasing in \(r\) at every strength (Proposition~\ref{prop:returns}(a)), \(\chi_r=1\) exactly when \(r\le r_C\). On
\((r_S,h)\), the posteriors after a sale and after no order are \(m<1/2<\tau(r)\), so entry after them is \(\rho\), and
when an indifferent high-cost challenger prepares, entry after a purchase is \(\rho+(1-\rho)\chi_r\), because
\(B_r(M)\ge c_H\) exactly when \(M\ge\tau(r)\). Weighting by the frequencies gives
\(\rho+(1-\rho)\chi_r\,(b+\lambda/2)\). At \(r_0\), Corollary~\ref{prop:region} gives a quiet market with entry
\(\rho\); at \(r_1\), \(\chi_{r_1}=1\). With a hidden price, entry is \(\rho\) after every order by
Proposition~\ref{prop:spread}(iv), because \(1/2<\tau(r)\).
\item The entry gain against \(r\) is \(w_r(\mu)=t_L-t_0+\DT(r)\mu\), with \(t_L-t_0=\ell-p-(\ell^2-3p^2)/(2r)\), which is
\(1/2-1/(8r)\) at the example's values. Since \(\ell^2-3p^2>0\) and \(\DT\) rises strictly above \(\ell\), both terms rise
strictly. A path of strengths in \((\ell,h)\) has lowest challenger value \(\ell\) and lowest strength above \(\ell\),
so \(\ell^2-3p^2>0\) is the condition \(p<\min\{\underline\theta,\underline r\}/\sqrt3\) of
Proposition~\ref{prop:spread}(vi)(b). At the prior \(1/2\), \(\bar f(1)=\bar f(-1)\), so
Proposition~\ref{prop:spread}(ii), with \(1/2<\tau(r)\), gives \(I_+-I_-=[e(1)-\rho]\,w_r(M)=(1-\rho)\chi_r\,w_r(M)\) by
(b), and \(I_+>I_-\) exactly when \(\chi_r=1\). On \((r_S,h)\) full two-sided trade holds, \(1/2<\tau(r)\), \(\DT\)
rises strictly and \(w_r(M)\) rises strictly, so Proposition~\ref{prop:spread}(vi)(c) applies on \((r_S,r_C]\), where
\(\chi_r=1\), and on \((r_C,h)\), where \(\chi_r=0\). By Proposition~\ref{prop:spread}(iv) and (ii), at \(r_1\) the
class component is \(\rho\DT(r_1)\vartheta\), the entry component is \((1-\rho)\,w_{r_1}(M)\), and the weighted impact
difference is \((1-\rho)(b+\lambda/2)\,w_{r_1}(M)\); with a hidden price the price gap is the class component. The
uniform laws have \(F_r(y)=\min\{y/r,1\}\), continuous in \(r\), and \((r_S,h)\) is a neighborhood of \(r_C\) on which
full two-sided trade and \(1/2<\tau(r)\) hold, so Proposition~\ref{prop:spread}(vi)(d) gives the fall
\((1-\rho)\,w_{r_C}(M)\), from \(\rho\DT(r_C)\vartheta+(1-\rho)\,w_{r_C}(M)\) to \(\rho\DT(r_C)\vartheta\). At \(r_2>r_C\),
\(\chi_{r_2}=0\). When an indifferent high-cost challenger may prepare with any probability, the last part of
Proposition~\ref{prop:spread}(vi)(d) applies, with an indifferent investor who may mix. Substituting the example's slot
in these expressions gives the values stated in (c).
\item Since \(\tau(r)>1/2\) at \(r_1\) and \(r_2\), Proposition~\ref{prop:informedshare}(b) applies at both. It gives the
left end \(3-2/v\) from the counts, and the right ends \(1-2k/(\rho\DT(r_1))\) and \(1-2k/(\rho\DT(r_2))\) from full
two-sided trade at each strength, both included when an indifferent investor may mix and the right end excluded when an
indifferent investor abstains. At the prior \(1/2\), a slot with share \(\vartheta'\) has \(M'=(1+\vartheta')/2\). Strict
recruitment at \(r_1\) is \(M'>\tau(r_1)\), and strict exclusion at \(r_2\) is \(M'<\tau(r_2)\); together they are the band.
Its left end exceeds \(3-2/v\): the inequality \(2\tau(r_1)-1>3-2/v\) is \(v<2/[4-2\tau(r_1)]=\reg{vol_band_max_frac}\),
and the example's volume lies below this bound by \(\reg{margin_vol_band_B}\). So
Proposition~\ref{prop:informedshare}(b) gives each share in the band a slot with the example's counts, and its entry
formula gives entry \(\rho+(1-\rho)v/2\) at \(r_1\) and \(\rho\) at \(r_2\), as at the example's slot, whose top
posterior lies in the band too. Such a slot has \(m'=(1-\vartheta')/2>1-\tau(r_2)\), so at \(r_1\) and at \(r_2\)
\[
\rho\,m'\,\DT(r)-k>\rho\,[1-\tau(r_2)]\,\DT(r_1)-k=\reg{margin_obs_band_incentive_B}>0,
\]
because \(\DT(r_2)>\DT(r_1)\). Proposition~\ref{prop:spread}(iv) then gives full two-sided trade as the unique outcome
of the slot at both strengths. The share \(\reg{alt_theta_B}\) lies in the band, and by
Proposition~\ref{prop:spread}(iv) its price gap at \(r_1\) is \(\rho\DT(r_1)\vartheta'+(1-\rho)\,w_{r_1}(M')\). On
\((r_S,h)\), entry after a sale and after no order is \(\rho\), and Proposition~\ref{prop:informedshare}(c) at the prior
\(1/2\) gives the share from the sale impact.
\item Corollary~\ref{cor:example-coexistence} and Table~\ref{tab:a4-set} give the two equilibria and their
intervals. At a full purchase with class \(L\) abstaining, both classes sell with probability \(b\), so a sale meets the
posterior \(1/2\) and earns \(\rho\DT(r)/2-k\), which is nonpositive exactly when \(r\le r_N\); for \(r<r_C\) the purchase
recruits and earns \(\DT(r)(1-M)-k\), which is nonnegative exactly when \(r\ge r_F\). On \([r_F,r_N]\) class \(H\) buys
with probability one, and Corollary~\ref{cor:purchasesonly}(a) at the prior
\(1/2\) gives \(\lambda=2\E[X]\). On \((r_A,r_F)\) class \(H\) mixes, an indifferent investor may mix, and the posterior
after a purchase exceeds \(\tau(r)\), so every challenger prepares after a purchase. Corollary~\ref{cor:purchasesonly}(b)
gives \(s_+=2b\,[\DT(r)/(2k)-1]=b\,[\DT(r)/k-2]\) and the failure to identify \(\lambda\).
\item For \(\ell\le r\le h\), \(R\le r\le h\), so a prepared high-class challenger wins with probability one, a prepared
low-class one wins exactly when \(R<\ell\), and \(W_L=\ell/r\). Proposition~\ref{prop:winner}(i) gives the win probability
after each order, and (ii) gives the winner gap \((M-m)\,\Omega=\vartheta(1-\ell/r)\) under full two-sided trade. Since
\(W_H=1\) and \(W_L=\ell/r\) falls strictly, Proposition~\ref{prop:winner}(v) makes the gap rise strictly. Here
\(D(y)=\1\{\ell\le y<h\}\), so \(D(r)-\Omega=\ell/r>0\) for \(\ell\le r<h\), the condition of
Proposition~\ref{prop:winner}(vi). On \((r_S,h)\) full two-sided trade holds with \(1/2<\tau(r)\) and \(\E[X]=0\), so
Proposition~\ref{prop:winner}(iii) and (iv) give \(\Cov(X,\1_{\rm prep})=(1-\rho)\chi_r(b+\lambda/2)\),
\(\Cov(X,\1_{\rm acq})=\rho\lambda\,\Omega/2+(1-\rho)\chi_r\,(b+\lambda/2)\,W_{F_r}(M)\), and the second difference
\((1-\rho)\chi_r\,W_{F_r}(M)\), with \(\chi_{r_1}=1\) and \(\chi_{r_2}=0\). Substituting the example's slot gives the
values stated in (f).
\item The incumbent laws strengthen with \(r\), so \(B_r(1/2)\) falls in \(r\) by Proposition~\ref{prop:returns}(a),
and \(B_r(1/2)<c_H\) for \(r\ge r_0\). Proposition~\ref{prop:announcement}(iii) gives a positive mean in every
equilibrium at each such strength. At \(r_1\) entry after the orders \((-1,0,1)\) is \((\rho,\rho,1)\), and at \(r_2\) it is
\(\rho\) after every order; the display of Proposition~\ref{prop:announcement}(iii) gives the two means. By
Proposition~\ref{prop:announcement}(iv) and Proposition~\ref{prop:spread}(ii), the expected price minus the price after
no order is \((1-\rho)\chi_r\,(b+\lambda/2)\,w_r(M)\), the weighted impact difference of (c), which is zero at \(r_2\).
Substituting the example's slot in \(P(x)=t_0+e(x)\,w_r(\mu(x))\) gives the prices stated in (g).
\item Lemma~\ref{lem:a4-order} gives \(r_N<r_C\), so \(r<r_C\). Under the default tie rules,
Corollary~\ref{cor:example-coexistence} and Table~\ref{tab:a4-set} give exactly two equilibria on \((r_F,r_N)\): a quiet
market and a full purchase with class \(L\) abstaining. In the second, Corollary~\ref{cor:purchasesonly}(a) at the prior
\(1/2\) with \(s_+=\lambda\) gives \(\bar f(1)\), \(\bar f(-1)\) and \(\E[X]=\lambda/2\). Both classes sell with
probability \(b\), so the posterior after a sale is \(1/2\). After no order, the class probabilities are \(1-2b-\lambda\)
and \(1-2b\), which give \(\mu(0)\) and
\(1/2-\mu(0)=\lambda/[2(2-4b-\lambda)]=\lambda/[4(1-v)]\). Both posteriors lie below \(\tau(r)\), so entry after both
orders is \(\rho\), and \(P(-1)-P(0)=\rho\,[w_r(1/2)-w_r(\mu(0))]=\rho\DT(r)\,[1/2-\mu(0)]\). After a purchase the
posterior is \(M>\tau(r)\), so every challenger prepares. In the quiet market every posterior is \(1/2\), entry is \(\rho\)
after every order, and the three prices coincide. When an indifferent high-cost challenger may prepare with any
probability, the threshold pair of Appendix~\ref{app:coexistence} is an equilibrium on \((r_A,r_J]\), which contains
\([r_F,r_N]\) by Lemma~\ref{lem:a4-order}. In it class \(L\) abstains; class \(H\) buys at the rate
\(s_+=b\,[2\tau(r)-1]/[1-\tau(r)]\), at which the posterior after a purchase equals \(\tau(r)\); and entry after a
purchase is \(e(1)=k/J(r)\), with \(J(r)=\DT(r)\,[1-\tau(r)]\), at which class \(H\) is indifferent. So class \(H\) buys
with probability \(s_+/\lambda\), the mean order is \(s_+/2\) by Corollary~\ref{cor:purchasesonly}(a), the low-class sale
earns \(\rho\DT(r)/2-k<0\), because the posterior after a sale is \(1/2\) and \(r<r_N\), and entry is
\(\rho+[e(1)-\rho]\,(b+s_+/2)\). Substituting the example's slot at \(r=\reg{str_r160}\) gives the values stated in (h).\qedhere
\end{enumerate}
\end{proof}
